\section{Discussion}
\label{sec:discussion}

To summarize, we compared research ideas generated by our AI agent with ideas written by expert researchers, and observed the robust finding that expert reviewers rate AI ideas as statistically more novel than expert ideas. In this section, we discuss some high-level questions readers might have and suggest some ways to address them. 


\textbf{Question 1: Do these collected expert ideas represent their best ideas?}
% \textbf{Compare AI ideas with published papers rather than new ideas brainstormed on the spot.}
One might argue that these ideas submitted by our idea-writing participants might not represent their best ideas as we discussed in Section 6.1, since most of them came up with the idea on the spot within a short period. 
In order to address this concern, we have designed an experiment where we will compare AI ideas with papers accepted at top-tier AI conferences. 
To avoid any possible contamination, we target the upcoming EMNLP 2024 conference, which will release the accepted papers in October 2024. 
We have generated AI ideas with our agent on 23 topics from the EMNLP Call For Papers page in July 2024 and cached them. 
We pre-registered our analysis plan which also includes the link to the cached ideas.~\footnote{\url{https://osf.io/z6qa4}}
Apart from comparing the quality of these ideas, we will also compute the overlap between AI-generated ideas and accepted papers on the same topics. 

\textbf{Question 2: Are evaluations based solely on ideas subjective?}
In this current study, we focused solely on evaluating the ideas themselves. Ideas that sound novel and exciting might not necessarily turn into successful projects, and our results indeed indicated some feasibility trade-offs of AI ideas. 
% Although we asked reviewers to assess the expected effectiveness of the ideas, these evaluations were inherently speculative, as reflected by the relatively low inter-annotator agreement we observed.
We view the current study as a preliminary evaluation of AI-generated ideas. In the next phase, we will recruit researchers to execute some AI and human-generated ideas into full projects. This will enable reviewers to assess the complete experimental outcomes, providing a more reliable basis for evaluation. Furthermore, it will allow us to analyze whether our initial idea evaluations align with the assessments of the actual project outcomes.
% \textbf{Extend from evaluating ideas to evaluating projects.} 
% Phase I of the study only evaluated the ideas themselves.  
% Ideas that sound novel and exciting don't necessarily turn into successful projects, and our results indicated some feasibility trade-offs of AI ideas. 
% Despite that we have asked reviewers to predict the expected effectiveness of the ideas as part of their evaluation, we acknowledge such evaluation tends to be speculative and we indeed observed a relatively low inter-annotator agreement. 
% We consider this phase I as a sanity check of the AI-generated ideas, and we will launch phase II of the study where we will recruit researchers to execute some of our AI and human ideas into full projects. 
% This way, reviewers can see the full set of experiment outcomes to make a more reliable judgment. 
% This will also allow us to retrospectively examine whether the evaluation results on the ideas themselves correlate with the evaluations based on the executed results. 


\textbf{Question 3: Why do you focus only  on prompting-based research in NLP?} 
The scope of our study is limited to prompting research ideas within NLP. We chose this design to facilitate the next phase of our execution experiment, where we prefer research ideas that are less resource-demanding and can be executed relatively quickly. 
We believe that the evaluation protocols we established should be applicable to other research domains as well, although the conclusions could be different depending on the research fields. 
Future work should consider extending such human study to other research domains and it would be interesting to compare how the conclusions differ. 

\textbf{Question 4: Can you automate idea execution as well?}
It is tempting to envision an end-to-end automated research pipeline where AI agents can implement AI-generated ideas to directly evaluate their effectiveness. 
Apart from speeding up scientific discovery, one could also imagine using such execution agents to automatically verify experiment results in existing papers or new submissions. 
We have also explored building an LLM agent to generate code to implement the generated ideas. 
Specifically, we provide a template codebase that consists of: (1) loading datasets from Huggingface or generating synthetic test examples; (2) implementing baseline methods; (3) implementing the proposed method; (3) loading or implementing the evaluation metrics; (4) running experiments on the testset with the baselines and the proposed method, so that
the output of the agent will be a report of the baseline performance as well as the proposed method's performance. 
While this agent can generate code that compiles and executes, we find that the automated experiments can be \textbf{misleading} because the agent often skips or modifies steps in the baselines or proposed methods. In some cases, the metric functions are also not correctly defined. 
This highlights the core challenge: just comparing the final experiment results is not enough; we have to verify the faithfulness of the implementations as well. 
Performing such implementation verification is not a trivial task, and we leave it to future work. 
We provide detailed description of our idea execution agent in Appendix~\ref{sec:failed_attempts}. 

% \textbf{Failed attempt on self-revision.} We also explored improving the novelty of generated ideas by asking the agent to compare the idea with top retrieved papers and revise the idea in order to make it more different. Such iterative revision leads to ideas that stack many different incremental changes together. While one may argue that such combinations could be novel, we decided not to adopt such iterative revision because we prefer one clearly novel idea over combinations of many known techniques.  







\section{Ethical Considerations}
\label{sec:ethics}


\textbf{Publication Policy.}
The growing use of AI to generate research ideas raises serious concerns about the potential abuse of these technologies by students or researchers who may flood academic conferences with low-quality or poorly thought-out submissions. The availability of LLM-generated content could lead to a decline in the overall quality of academic discourse, as some individuals might take a lazy approach, relying on AI to both generate ideas and review submissions. This would undermine the credibility and integrity of the review process.
%
The risks are real. Without proper oversight, we could see a deluge of submissions that lack depth or intellectual merit. To prevent this, it is essential to hold researchers accountable for the outputs generated through AI tools. Rigorous standards must be applied equally to both AI-assisted and human-generated research to ensure that the use of LLMs does not result in misleading, superficial, or unethical academic contributions.
% The review and publication of AI-generated research ideas introduce new challenges regarding the credibility, transparency, and ethical integrity of academic publications. 
% Questions would also arise about how to attribute authorship and whether AI should be listed as a co-author or simply acknowledged as a tool.
% While it is open for debate whether ideas originated from AI should be published at the same venues as human ideas, we believe the bottom line is that we should carefully examine all the research outputs with the same level of rigor and hold the researchers accountable if their use of LLMs in the research process leads to misleading results. 




\textbf{Intellectual Credit.}
The use of LLMs to generate research ideas introduces significant ambiguity around the concept of intellectual credit. Traditional frameworks for attributing credit in research, based on human authorship and contribution, become less clear when AI plays a significant role in idea generation. Questions arise around how to distribute credit between the developers of the LLM, the researchers who designed the frameworks for its use, and the researchers who integrate AI-generated ideas into their work.
%
Furthermore, it becomes increasingly difficult to trace the origins of AI-generated contributions, especially when they draw from vast datasets composed of numerous sources. This complexity calls for a broader rethinking of how intellectual credit is assigned in AI-driven research.
%
While a complete overhaul of legal and academic norms is beyond the scope of this project, we advocate for the adoption of transparent documentation practices. Researchers should clearly disclose the role AI played in the idea generation process, specifying which models, data sources, and frameworks were used, and outlining the level of human involvement. This could ensure that the credit distribution in AI-supported research is as transparent and fair as possible.


\textbf{Potential for Misuse.}
AI-generated research ideas, especially those that introduce novel concepts, have the potential to be misused in ways that could lead to harmful or destabilizing outcomes. For instance, ideation agents could be leveraged to generate adversarial attack strategies or other unethical applications. This concern aligns with broader arguments from those focused on existential risk (X-risk), who argue that AI-driven innovation could be a primary route to destabilizing the status quo, posing risks at a societal or even global level.
%
Our stance is that such discussions on safety should be evidence-based to the extent that it is possible, and careful evaluation work is an important component of keeping these discussions grounded in actual, measured capabilities of these systems. We advocate for continued safety research specifically targeting these types of concerns—such as the development of Reinforcement Learning from Human Feedback (RLHF) systems or anti-jailbreak mechanisms for research ideation agents. Additionally, we believe it would be meaningful to create safety benchmarks that assess the ethical and safe application of AI-generated ideas.


\textbf{Idea Homogenization.}
Our analysis showed that current LLMs lack diversity in idea generation. 
This raises important concerns that wide adoption of LLMs 
 can result in idea homogenization, where the generated ideas only reflect a narrow set of perspectives or have systematic biases. 
 Over time, this could lead to a reduction in the richness and diversity of research outputs globally.
 Future work should develop ways to either improve LLMs themselves or refine our idea generation methods to promote idea diversity. 
 It's also important to note that our evaluation primarily assesses the quality of the typical ideas being generated, and may not fully capture the long tail of unique or novel ideas that would be truly transformative.
 % We also highlight the importance of including diversity as part of the evaluation for all future work. 


\textbf{Impact on Human Researchers.}
The integration of AI into research idea generation introduces a complex sociotechnical challenge, as research is fundamentally a community-driven, collaborative effort. By introducing AI, particularly LLMs, into this social system, we risk unforeseen consequences. 
% For instance, there is a potential for displacement of human researchers, the devaluation of human creativity, and the disruption of collaborative research dynamics. 
Overreliance on AI could lead to a decline in original human thought, while the increasing use of LLMs for ideation might reduce opportunities for human collaboration, which is essential for refining and expanding ideas.
%
To mitigate these risks, future works should explore new forms of human-AI collaboration, and our results on human reranking of AI ideas show that even naive human-AI collaboration approaches can be effective. Beyond reranking, humans can play a critical role in the ideation process by providing intermediate feedback, taking AI-generated ideas as inspiration for further development, and bringing their unique expertise into the process. %Importantly, human collaboration in research often involves brainstorming and refining ideas together—an element that AI alone cannot replicate.
%
Understanding how to integrate LLMs into this collaborative process without disrupting the social fabric of research will be an important ongoing problem, requiring careful consideration of the broader sociotechnical implications.


% \textbf{Impact on Human Researchers.}
% The use of AI to generate research ideas raises concerns about the potential displacement of human researchers and the devaluation of human creativity. There is a risk that researchers may become overly reliant on AI, leading to a decline in original human thought and innovation. Furthermore, the dynamics of research collaboration could be fundamentally altered. For example, increasing use of LLMs for ideation might discourage collaboration among human researchers. 
% To address this, we highlight the value of human-AI collaboration. We presented preliminary results where human reranking on top of AI-generated ideas can bring additional values. 
% Apart from reranking, there are many other possible ways for humans to contribute to the collaborative ideation process, for example, by providing intermediate feedback to generated ideas, or taking AI ideas as inspirations for further improvement. 
% Moreover, human researchers often brainstorm together and collaborative discussion helps refine ideas. 
% How to adapt LLMs in collaborative idea generation is an interesting open question that we leave to future work. 





\section*{Positionality Statement}

We disclose the authors' anticipated outcomes of the human study before the experiment was conducted to be transparent about experimenter biases. 
Among the three authors, Tatsu and Diyi were expecting a null result from the study while Chenglei expected AI to be better than humans. 
% positions on AI to be transparent about any potential biases. 
%In terms of the authors' general stances on AI, Chenglei believes in AGI, Tatsu loves scaling, and Diyi prioritizes being human-centered. 